\chapter{Guia complementar de adaptação}\label{ap:complementar}

\section{Elicitação, perfis e jornada}\label{guia:elicitacao}
Uma SRS liga uma necessidade a uma obrigação verificável. Ao adaptar o modelo,
registre problema, evidência, impacto e indicador de sucesso antes de escolher
a solução. Diferencie evidência observada, hipótese e decisão do produto.
Uma origem de requisito identifica pessoa, documento ou regra; não demonstra,
por si só, que uma entrevista ocorreu.

\begin{quadrosrs}
\caption{Técnicas de elicitação e registros úteis}
\begin{tabelasrs}{L{0.22\textwidth}Y}
\cabecalhosrs\textbf{Técnica} & \textbf{Registro para análise e validação}\\\midrule
Entrevista & Papel do participante, roteiro, data, consentimento aplicável e síntese validada.\\
Observação & Contexto, tarefa, dificuldade observada e diferença entre relato e comportamento.\\
Oficina & Participantes, decisões, conflitos e fonte de cada regra acordada.\\
Documentos e dados & Versão, procedência, limites de interpretação e regras extraídas.\\
Protótipo & Tarefa avaliada, perfil do usuário, resultado observado e alteração justificada.\\
\end{tabelasrs}
\fonte{Elaboração didática deste guia.}
\end{quadrosrs}

Descreva perfis por objetivo, contexto, frequência de uso, restrição de acesso e
necessidade de acessibilidade. Uma persona resume essas características com a
sua fonte; não substitui uma pessoa real na validação. Para cada stakeholder,
registre representante, interesse, responsabilidade, influência e modo de
participação, além dos papéis exemplificados na \autoref{cap:descricao}.

Uma jornada pode organizar etapa, ação da pessoa, resposta do sistema,
dificuldade e oportunidade. Neste domínio, consultar o catálogo, confirmar a
inscrição, consultar seus registros e cancelar correspondem a \ucref{UC-003},
\ucref{UC-001}, \ucref{UC-004} e \ucref{UC-002}. Cada etapa deve apontar para
requisitos e aceite, evitando criar funcionalidades apenas no desenho.

Uma história exemplificativa é: ``Como participante, quero consultar minhas
inscrições para acompanhar sua situação''. Sua obrigação está em \reqref{RF-005}
e sua verificação em \ctref{CT-008}. Histórias descrevem valor; as fichas e os
cenários Dado/Quando/Então delimitam comportamento e exceções.

\clearpage
\section{Qualidade dos requisitos e da documentação}\label{guia:qualidade}
Revise cada enunciado quanto a atomicidade, clareza, viabilidade, consistência,
necessidade, rastreabilidade e verificação. Condições e unidades eliminam
interpretações como ``sempre disponível'' ou ``tempo aceitável''. Além das
categorias demonstradas pelos RNF deste modelo, avalie as necessidades do produto:

\begin{quadrosrs}
\caption{Categorias complementares de qualidade}
\begin{tabelasrs}{L{0.24\textwidth}Y}
\cabecalhosrs\textbf{Categoria} & \textbf{Pergunta para especificar e medir}\\\midrule
Usabilidade e acesso & Quais tarefas, perfis, tecnologias assistivas, contraste, tamanho de texto e treinamento precisam de verificação?\\
Confiabilidade & Qual falha é medida? Qual janela, taxa de erro, disponibilidade, precisão de resultado ou tempo de reparo é exigido?\\
Desempenho & Qual carga, capacidade, taxa de transferência, uso de recursos, percentil e janela delimitam o resultado?\\
Suportabilidade & Que diagnóstico, configuração, instalação, evolução e documentação permitem manter o produto?\\
Interoperabilidade & Quais versões, formatos, limites, dependências e falhas de integração precisam de contrato?\\
Restrições de design & Qual plataforma, padrão ou condição de implantação é obrigatória e qual é a origem dessa restrição?\\
\end{tabelasrs}
\fonte{Elaboração didática deste guia.}
\end{quadrosrs}

Disponibilidade, tempo médio entre falhas e tempo médio de reparo exigem
definições da população, do período e da falha. Não podem ser deduzidos apenas
da escolha de Docker, de uma rede local ou de um fornecedor. Este guia não
adiciona metas dessas categorias ao exemplo: uma nova obrigação exige ID,
origem, métrica, meta e cenário de verificação.

Especifique o comportamento em modo degradado e como identificar seu início e
término. Para manutenção, relacione convenções de código, nomes, bibliotecas e
utilitários à necessidade de suporte e à sua verificação.

Documentação do produto pode incluir ajuda contextual, manual de uso,
instalação, operação e contrato de API. Para cada artefato, defina público,
conteúdo, formato, idioma, acesso, responsabilidade e critério de atualização.
Defina também os resultados de consultas ou relatórios: campos e origem,
filtros, ordenação, período, fuso, permissão, exportação e aceite. No exemplo,
a representação de inscrições pessoais já é delimitada por \reqref{RF-005}.

Na adaptação, identifique licença do software e dos recursos de terceiros,
direitos de uso de conteúdo e marcas e normas aplicáveis com edição e cláusula.
Relacione cada obrigação a uma fonte e à sua verificação. A licença do modelo
documental não define a licença da aplicação que ele descreve. Usar uma classe
de formatação também não demonstra conformidade integral com uma norma.

\clearpage
\section{Arquitetura, dados e decisões}\label{guia:arquitetura}
Escolha modelos de acordo com a pergunta. Uma visão de contexto mostra
fronteiras; uma arquitetura funcional mostra capacidades; componentes e
camadas mostram responsabilidades; classes e dados mostram conceitos e
relações. Uma vista C4, quando utilizada, é complementar à UML. O nível de
detalhe deve permitir conferir a correspondência com as fichas.

\begin{quadrosrs}
\caption{Obrigações que orientam decisões de design}
\begin{tabelasrs}{L{0.25\textwidth}Y}
\cabecalhosrs\textbf{Origem} & \textbf{Questão de arquitetura ou dados}\\\midrule
\rnref{RN-001} & Como garantir capacidade e unicidade ativa na mesma operação, inclusive sob concorrência?\\
\reqref{RNF-001} & Como verificar latência nas condições de carga e recursos especificadas?\\
\reqref{RNF-002} & Em qual fronteira são verificadas sessão e titularidade em toda operação pessoal?\\
\reqref{RNF-004} & Como separar ciclo de vida do processo e conservação dos dados confirmados?\\
\end{tabelasrs}
\fonte{Elaboração didática deste guia.}
\end{quadrosrs}

O modelo conceitual descreve o domínio. O modelo lógico define entidades,
chaves, cardinalidades e nulabilidade. O modelo físico escolhe tipos do banco,
índices, restrições e migrações. A \autoref{fig:dados-logicos} e o dicionário
são base para esse mapeamento; não equivalem a um script SQL executável.

Neste exemplo, \texttt{eventoId} pode referenciar um evento local, mas
\texttt{participanteId} identifica um titular externo. A implementação deve
garantir no máximo uma inscrição \emph{ativa} por titular e evento, preservando
os registros cancelados. Uma restrição de unicidade sobre todo o histórico
impediria a reinscrição permitida por \reqref{RF-004}. Tipos, unidades,
sensibilidade e regras de exclusão precisam acompanhar cada campo no esquema
físico. Não transportar automaticamente cascatas, senhas ou tabelas de outro
domínio para essa aplicação.

Registre decisões arquiteturais em ADRs: título, contexto e forças, opções
consideradas, decisão, estado da decisão e consequências. Evidências e
requisitos justificam a escolha. Um ADR trata uma decisão; alterações de
contrato precisam atualizar a SRS e seus testes. A política de versionamento
de dados, migrações e reversão também deve relacionar esquema e versão do serviço.

\clearpage
\section{Segurança e privacidade}\label{guia:seguranca}
Classifique os dados e delimite acesso, finalidade, retenção e descarte. A
classificação informa o design e os testes; não substitui o contrato de
titularidade da \reqref{RNF-002}.

\begin{quadrosrs}
\caption{Classificação exemplificativa dos dados}
\begin{tabelasrs}{L{0.19\textwidth}YY}
\cabecalhosrs\textbf{Classe} & \textbf{Exemplo} & \textbf{Cuidado de especificação}\\\midrule
Público & Eventos publicados e vagas. & Delimitar campos e momento de publicação.\\
Interno & Configuração e diagnóstico. & Definir acesso e retirar detalhes internos de erros públicos.\\
Pessoal & Titular e suas inscrições. & Minimizar coleta e aplicar acesso do próprio titular.\\
Segredo & Credencial de sessão ou integração. & Proteger uso e evitar gravação em logs ou documentos.\\
\end{tabelasrs}
\fonte{Elaboração didática deste guia.}
\end{quadrosrs}

Para cada dado pessoal, registre finalidade, necessidade, fundamento aplicável,
prazo ou evento de retenção, responsáveis, autorização de acesso e procedimentos
de consulta, exportação e exclusão. Delimite também dados de telemetria e
terceiros que recebem informações. O domínio usa identidade externa e não
especifica cadastro de senha local. A conservação do histórico no exemplo não
estabelece retenção ilimitada nem inventa uma política legal de descarte.

Analise ameaças por caminho de acesso: identidade falsa, alteração indevida,
negação de uma operação realizada, divulgação de informações, esgotamento de
recursos e elevação de privilégio. Registre ativo, atacante ou condição,
consequência, controle e cenário de verificação. Um catálogo de controles pode
ajudar a revisão, mas a seleção precisa estar ligada ao contexto e a uma versão
identificada; mencionar um catálogo não comprova atendimento.

No exemplo, \ctref{CT-006} verifica sessão, titularidade e rejeição de titular
forjado no corpo. \ctref{CT-002} verifica concorrência e capacidade.
Controles de transporte, limites de requisição, proteção de segredos e auditoria
que forem exigidos na adaptação devem receber seus próprios critérios de aceite.
Não deduza uma certificação de segurança ou acessibilidade a partir desses testes.

\clearpage
\section{Riscos e priorização}\label{guia:riscos}
Um registro de risco relaciona condição, probabilidade, impacto, resposta,
responsável e gatilho de acompanhamento. As situações abaixo ilustram como
ligar risco a obrigações já existentes; não representam uma avaliação realizada
com uma equipe ou dados de produção.

\begin{fichasrs}{Exemplos de riscos ligados à especificação}
\campoficha{Última vaga}{Condição: duas pessoas confirmam simultaneamente. Impacto: exceder a capacidade. Resposta: atomicidade e verificação de \rnref{RN-001} por \ctref{CT-002}. Responsabilidade de acompanhamento: equipe do produto. Gatilho: divergência entre capacidade e total ativo.}
\campoficha{Acesso indevido}{Condição: tentativa de consultar ID de outra pessoa. Impacto: exposição de registro pessoal. Resposta: sessão e titularidade em toda operação; \ctref{CT-006}. Responsabilidade de acompanhamento: equipe do produto. Gatilho: resposta pessoal sem autorização.}
\campoficha{Perda confirmada}{Condição: reinício normal do serviço. Impacto: sumiço de um registro confirmado. Resposta: persistência independente do processo e \ctref{CT-009}. Responsabilidade de acompanhamento: operação do produto. Gatilho: diferença entre registros antes e depois do reinício.}
\campoficha{Avaliação}{Para um projeto concreto, atribuir ID, origem, probabilidade e impacto com escala e justificativa, resposta, responsável nominal, data de análise e condição de reavaliação.}
\end{fichasrs}

A escala Alta/Média/Baixa da \autoref{cap:introducao} permanece como prioridade
das fichas. Outras técnicas podem apoiar a seleção de uma entrega: separar
essencial, desejável e opcional; comparar valor e esforço; considerar risco,
urgência e dependências. Quando usar pontuação, declare fórmula, unidades e
fontes. Uma pontuação não altera automaticamente uma obrigação acordada.

Mudanças de escopo precisam de análise de impacto sobre dados, interfaces,
modelos, testes e custo. Registre a decisão e mantenha os IDs, inclusive quando
uma obrigação for substituída. A matriz completa permanece na
\autoref{tab:rastreabilidade}; um diagrama ilustra relações, mas não dispensa
a conferência de cobertura.

\clearpage
\section{Processo, cronograma e custos}\label{guia:planejamento}
Escolha o processo conforme incerteza, feedback, tamanho da equipe e exigência
de aprovação. Entregas incrementais permitem validar partes; Kanban organiza
fluxo e limites de trabalho; Scrum organiza ciclos e responsabilidades;
um processo sequencial organiza etapas; uma abordagem híbrida combina práticas
com suas justificativas. A escolha não decorre apenas do nome da disciplina.

\begin{quadrosrs}
\caption{Campos de planejamento e estimativa}
\begin{tabelasrs}{L{0.24\textwidth}Y}
\cabecalhosrs\textbf{Artefato} & \textbf{Conteúdo necessário}\\\midrule
Mapa de entregas & Objetivo, escopo por IDs, dependências e critério de saída de cada incremento.\\
Cronograma & Atividade, início, término, dependência, responsável e marco verificável; um Gantt representa essas relações.\\
Estimativa & Unidade, intervalo, premissas, método, capacidade disponível e fonte da estimativa.\\
Responsabilidades & Quem executa, quem responde pela decisão, quem é consultado e quem recebe informação.\\
Custos & Quantidade, unidade, taxa, período, fonte e subtotal de cada categoria.\\
\end{tabelasrs}
\fonte{Elaboração didática deste guia.}
\end{quadrosrs}

Uma matriz RACI pode registrar responsável pela execução (R), responsável final
(A), consultados (C) e informados (I). Defina uma autoridade final por atividade
e confirme os papéis com as pessoas. Stakeholder, representante e responsável
pela tarefa não são necessariamente o mesmo papel.

Expresse esforço em horas ou pessoas-dia, deixando explícita a conversão e a
capacidade. Esforço agregado não é igual ao prazo de calendário. Considere
dependências, revisão, integração e testes ao estimar intervalos. Não apresente
datas, nomes ou valores de um exemplo como compromissos reais.

Calcule custo de pessoas como soma de horas por papel multiplicadas pela taxa
por hora. Some infraestrutura, ferramentas, serviços e outros custos do mesmo
período. Declare hipóteses de contingência e evite contar um custo duas vezes.
Apresente subtotais por etapa e total, moeda, período e fontes das taxas.
Estabeleça como uma alteração de escopo afeta essas estimativas.

\clearpage
\section{Configuração, integração e entrega}\label{guia:entrega}
Controle versões de código, especificação, diagramas, esquema e configuração
de infraestrutura. Identifique mudanças por data, versão, autor, descrição e
IDs afetados; o histórico Git permite recuperar versões. Credenciais e dados
pessoais de produção não pertencem ao repositório. Defina política de branches,
revisão por PR, versões, tags e notas de entrega conforme o projeto.

\begin{quadrosrs}
\caption{Ambientes e evidência de entrega da aplicação}
\begin{tabelasrs}{L{0.18\textwidth}YY}
\cabecalhosrs\textbf{Ambiente} & \textbf{Uso} & \textbf{Dados e evidência}\\\midrule
Local & Desenvolvimento e diagnóstico. & Dados sintéticos e configuração documentada.\\
Integração contínua & Build e verificações reproduzíveis. & Dados isolados, logs e resultados ligados à versão.\\
Homologação & Validação de contratos e fluxos integrados. & Dados controlados e diferenças de ambiente explícitas.\\
Produção & Operação do produto. & Acesso autorizado, métricas e procedimentos de recuperação.\\
\end{tabelasrs}
\fonte{Elaboração didática deste guia.}
\end{quadrosrs}

Um pipeline da aplicação pode encadear análise estática, testes unitários,
verificações de segurança, build, testes de integração e contrato, ambiente
temporário, testes de ponta a ponta e critérios de implantação. Registre
artefato e resultado de cada etapa. O pipeline deste repositório compila e
verifica o \emph{documento}; ele não executa os cenários do sistema fictício.

Cada integração precisa delimitar formato, versão, autenticação, erros,
limites, tempo de espera e política de repetição. Quando aplicável, descreva
TLS, portas, restrições de rede, hardware e comportamento sem conectividade.
Uma repetição precisa considerar idempotência, intervalo, limite e falha final.
O contrato HTTP existente define resultados e precedência; parâmetros adicionais
de operação devem ser acordados na adaptação, com verificação própria.

Associe versões da aplicação e do esquema às migrações. Defina implantação,
checagem de saúde, observação após a entrega e reversão. Rollback de código não
reverte necessariamente mudanças de dados. Teste a compatibilidade e descreva
como preservar registros confirmados durante uma recuperação.

\clearpage
\section{Observabilidade, suporte e recuperação}\label{guia:operacao}
Uma operação precisa de sinais que permitam identificar falhas e seu impacto.
Defina eventos, métricas e procedimentos com finalidade, responsável e acesso.

\begin{quadrosrs}
\caption{Sinais e procedimentos para especificar a operação}
\begin{tabelasrs}{L{0.24\textwidth}Y}
\cabecalhosrs\textbf{Item} & \textbf{Definição e verificação}\\\midrule
Erros e logs & Mensagem segura ao usuário, diagnóstico interno e identificador de correlação; excluir segredos e minimizar dados pessoais.\\
Métricas & Latência, taxa de erro, volume e saturação; filas e tarefas, se existirem; indicador de negócio com unidade e janela.\\
Rastreamento & Como relacionar uma solicitação às operações e integrações, com política de coleta e retenção.\\
Suporte & Canal, horário, classificação de incidente, resposta, escalonamento e comunicação de resultado.\\
Backup e restauração & Dados cobertos, frequência, proteção, retenção, procedimento e teste de restauração.\\
Recuperação & Perda tolerável de dados (RPO), tempo tolerável para recuperar (RTO), responsáveis e evidência de execução.\\
\end{tabelasrs}
\fonte{Elaboração didática deste guia.}
\end{quadrosrs}

Um indicador de nível de serviço (SLI) define o que é medido; um objetivo de
nível de serviço (SLO) define limite e janela. Disponibilidade, latência e erro
precisam indicar população, exclusões e método. Alertas devem apontar ação e
responsável. Não confunda percentil de resposta com disponibilidade nem com
uma garantia contratual de atendimento.

O ensaio de \reqref{RNF-001} tem carga e janela próprias; não estabelece um SLO
de produção. \reqref{RNF-004} cobre reinício normal com armazenamento preservado;
não demonstra recuperação de desastre ou perda do armazenamento. Metas de
backup, RPO e RTO seriam obrigações adicionais, a especificar e verificar
quando exigidas pelo contexto.

\clearpage
\section{Portfólio de testes, aceite e revisão}\label{guia:aceite}
Escolha testes pela obrigação e pelo tipo de falha. Registre ID do cenário,
requisito, objetivo, condições, dados, passos, resultado esperado, tipo,
responsável, localização da automação e evidência. Um resultado de execução
identifica versão, ambiente, data e conclusão; o catálogo de cenários da SRS
especifica o teste, sem afirmar que a aplicação foi executada.

\begin{quadrosrs}
\caption{Níveis de verificação e exemplos de finalidade}
\begin{tabelasrs}{L{0.22\textwidth}Y}
\cabecalhosrs\textbf{Nível ou foco} & \textbf{Falha que ajuda a identificar}\\\midrule
Unidade & Cálculo e regras de elegibilidade em condições de limite.\\
Integração & Atomicidade e conservação de registros reais no armazenamento.\\
Contrato & Formatos, erros, permissões e versões nas interfaces.\\
Ponta a ponta & Jornada completa entre interface e serviço com dados controlados.\\
Qualidade & Carga, acesso indevido, teclado e reinício nas condições dos RNF.\\
\end{tabelasrs}
\fonte{Elaboração didática deste guia.}
\end{quadrosrs}

Relacione obrigações a componentes responsáveis, cenários e evidências, além
de objetivo e origem. A matriz e os procedimentos da \autoref{cap:verificacao}
incluem alternativas, exceções, última vaga, repetição e limite exato de prazo.
Confirme que uma alteração dos modelos não remove um desses caminhos.

Critérios de conclusão e de aceite de uma entrega podem exigir revisão de
alterações, build e testes aprovados, requisitos cobertos, risco analisado,
segurança verificada, migração e reversão exercitadas, sinais de operação e
documentação acessível. Defina quem avalia e qual evidência demonstra cada
critério; uma lista preenchida sem evidência não comprova aprovação.

Uma revisão acadêmica pode avaliar problema e escopo, qualidade dos requisitos,
consistência dos modelos, adequação das decisões, qualidade e segurança,
verificação e operação. Para cada dimensão, declare critério, evidência e escala
de avaliação. Não transforme a rubrica em resultados atribuídos a este exemplo.

Organize os artefatos por finalidade: contexto e fontes na descoberta;
requisitos, contratos e aceite na especificação; modelos e ADRs no design;
código e migrações na implementação; resultados na verificação; versão,
procedimentos e indicadores na operação. Reavalie relações e vocabulário a cada
mudança de escopo. Referências com versão e procedência tornam essa revisão
rastreável.
